% GENERATED FILE. Edit canonical lessons, then run npm run generate:books.
% canonical-source-hash: fnv1a64:11d26d66ea68a079
% canonical-lessons: HI-C550-marzi, HI-C550-manpasand, HI-C550-manahi, HI-C550-pasandida, HI-C550-cayan

\chapter{Making Up One's Mind}
\label{ch:hi-making-up-one-s-mind}
\hlchaptermodality{550}

\begin{chapteropening}
\textbf{By the end of this chapter:} \emph{I can say what I chose, accepted or refused.}

Complete the last lesson of chapter 550: I can say what I chose, accepted or refused.
\end{chapteropening}

\section[\texorpdfstring{\hi{मर्ज़ी} (marzī)}{मर्ज़ी (marzī)}]{\hi{मर्ज़ी} (marzī) — one's own wish; one's choice}

\label{lesson:HI-C550-marzi}



\begin{quote}
Before the new one: say the Hindi for preference, then the Hindi for to hesitate, to hold back.
\end{quote}

\subsection*{You'll want to know: \hi{मर्ज़ी}}
\textbf{\hi{मर्ज़ी}} — \emph{marzī} — \textquotedblleft{}one's own wish; one's choice\textquotedblright{}.

\begin{quote}
\hi{छुट्टी} \hi{के} \hi{दिन} \hi{मैं} \hi{अपनी} \hi{मर्ज़ी} \hi{से} \hi{उठता} \hi{हूँ।}
\end{quote}

\begin{quote}
\emph{Chuṭṭī ke din maiṁ apnī marzī se uṭhtā hūṁ.}
\end{quote}

\begin{quote}
On a day off I get up when I like.
\end{quote}

Feminine.

\begin{cousinweb}
\textbf{\hi{मर्ज़ी}} (\emph{marzī}) is usually traced to Arabic \emph{marḍī}, \textquotedblleft{}pleasing, approved\textquotedblright{}, through Persian.
\end{cousinweb}

\begin{grammarlens}[title={What you've built}]
One more word for saying what happens if.
\end{grammarlens}

\subsection*{Your turn}
\begin{itemize}
\raggedright
  \item \emph{Say it:} \emph{marzī}
  \item \emph{Say it:} \emph{marzī}, once more
  \item \emph{Say it:} \emph{Chuṭṭī ke din maiṁ apnī marzī se uṭhtā hūṁ.} — On a day off I get up when I like.
  \item \emph{Recall:} say the Hindi for preference, then the Hindi for to hesitate, to hold back, then say \emph{marzī} again
\end{itemize}

\begin{tcolorbox}[breakable,colback=teal!4,colframe=teal!35!black,title={Before you move on}]
What does \hi{मर्ज़ी} mean? (\textquotedblleft{}One's own wish; one's choice\textquotedblright{}.) Say it once more.
\end{tcolorbox}
\section[\texorpdfstring{\hi{मनपसंद} (manpasand)}{मनपसंद (manpasand)}]{\hi{मनपसंद} (manpasand) — of one's own choosing, favourite}

\label{lesson:HI-C550-manpasand}



\begin{quote}
Before the new one: say the Hindi for to hesitate, to hold back, then the Hindi for one's own wish.
\end{quote}

\subsection*{You'll want to know: \hi{मनपसंद}}
\textbf{\hi{मनपसंद}} — \emph{manpasand} — \textquotedblleft{}of one's own choosing, favourite\textquotedblright{}.

\begin{quote}
\hi{मैंने} \hi{अपनी} \hi{मनपसंद} \hi{किताब} \hi{ख़रीदी।}
\end{quote}

\begin{quote}
\emph{Maiṁne apnī manpasand kitāb kharīdī.}
\end{quote}

\begin{quote}
I bought the book I wanted.
\end{quote}

\begin{cousinweb}
\textbf{\hi{मनपसंद}} (\emph{manpasand}) joins \emph{man}, \textquotedblleft{}mind\textquotedblright{}, from Sanskrit \emph{manas} (CDIAL 9822), and \emph{pasand}, \textquotedblleft{}liked\textquotedblright{}, which is usually traced to Persian \emph{pasand}: what the mind likes.
\end{cousinweb}

\begin{grammarlens}[title={What you've built}]
One more word for saying what happens if.
\end{grammarlens}

\subsection*{Your turn}
\begin{itemize}
\raggedright
  \item \emph{Say it:} \emph{manpasand}
  \item \emph{Say it:} \emph{manpasand}, once more
  \item \emph{Say it:} \emph{Maiṁne apnī manpasand kitāb kharīdī.} — I bought the book I wanted.
  \item \emph{Recall:} say the Hindi for to hesitate, to hold back, then the Hindi for one's own wish, then say \emph{manpasand} again
\end{itemize}

\begin{tcolorbox}[breakable,colback=teal!4,colframe=teal!35!black,title={Before you move on}]
What does \hi{मनपसंद} mean? (\textquotedblleft{}Of one's own choosing, favourite\textquotedblright{}.) Say it once more.
\end{tcolorbox}
\section[\texorpdfstring{\hi{मनाही} (manāhī)}{मनाही (manāhī)}]{\hi{मनाही} (manāhī) — a ban, a prohibition}

\label{lesson:HI-C550-manahi}



\begin{quote}
Before the new one: say the Hindi for one's own wish, then the Hindi for of one's own choosing.
\end{quote}

\subsection*{You'll want to know: \hi{मनाही}}
\textbf{\hi{मनाही}} — \emph{manāhī} — \textquotedblleft{}a ban, a prohibition\textquotedblright{}.

\begin{quote}
\hi{मंदिर} \hi{में} \hi{जूते} \hi{पहनकर} \hi{जाने} \hi{की} \hi{मनाही} \hi{है।}
\end{quote}

\begin{quote}
\emph{Mandir meṁ jūte pahankar jāne kī manāhī hai.}
\end{quote}

\begin{quote}
Wearing shoes inside the temple is not allowed.
\end{quote}

Feminine.

\begin{cousinweb}
\textbf{\hi{मनाही}} (\emph{manāhī}) is usually traced to Arabic \emph{manāhī}, \textquotedblleft{}forbidden things\textquotedblright{}, through Persian; Hindi uses it for a ban.
\end{cousinweb}

\begin{grammarlens}[title={What you've built}]
One more word for saying what happens if.
\end{grammarlens}

\subsection*{Your turn}
\begin{itemize}
\raggedright
  \item \emph{Say it:} \emph{manāhī}
  \item \emph{Say it:} \emph{manāhī}, once more
  \item \emph{Say it:} \emph{Mandir meṁ jūte pahankar jāne kī manāhī hai.} — Wearing shoes inside the temple is not allowed.
  \item \emph{Recall:} say the Hindi for one's own wish, then the Hindi for of one's own choosing, then say \emph{manāhī} again
\end{itemize}

\begin{tcolorbox}[breakable,colback=teal!4,colframe=teal!35!black,title={Before you move on}]
What does \hi{मनाही} mean? (\textquotedblleft{}A ban, a prohibition\textquotedblright{}.) Say it once more.
\end{tcolorbox}
\section[\texorpdfstring{\hi{पसंदीदा} (pasandīdā)}{पसंदीदा (pasandīdā)}]{\hi{पसंदीदा} (pasandīdā) — favourite}

\label{lesson:HI-C550-pasandida}



\begin{quote}
Before the new one: say the Hindi for of one's own choosing, then the Hindi for a ban, a prohibition.
\end{quote}

\subsection*{You'll want to know: \hi{पसंदीदा}}
\textbf{\hi{पसंदीदा}} — \emph{pasandīdā} — \textquotedblleft{}favourite\textquotedblright{}.

\begin{quote}
\hi{आम} \hi{मेरा} \hi{पसंदीदा} \hi{फल} \hi{है।}
\end{quote}

\begin{quote}
\emph{Ām merā pasandīdā phal hai.}
\end{quote}

\begin{quote}
The mango is my favourite fruit.
\end{quote}

\begin{cousinweb}
\textbf{\hi{पसंदीदा}} (\emph{pasandīdā}) is usually traced to Persian \emph{pasandīda}, \textquotedblleft{}approved, chosen\textquotedblright{}.
\end{cousinweb}

\begin{grammarlens}[title={What you've built}]
One more word for saying what happens if.
\end{grammarlens}

\subsection*{Your turn}
\begin{itemize}
\raggedright
  \item \emph{Say it:} \emph{pasandīdā}
  \item \emph{Say it:} \emph{pasandīdā}, once more
  \item \emph{Say it:} \emph{Ām merā pasandīdā phal hai.} — The mango is my favourite fruit.
  \item \emph{Recall:} say the Hindi for of one's own choosing, then the Hindi for a ban, a prohibition, then say \emph{pasandīdā} again
\end{itemize}

\begin{tcolorbox}[breakable,colback=teal!4,colframe=teal!35!black,title={Before you move on}]
What does \hi{पसंदीदा} mean? (\textquotedblleft{}Favourite\textquotedblright{}.) Say it once more.
\end{tcolorbox}
\section[\texorpdfstring{\hi{चयन} (cayan)}{चयन (cayan)}]{\hi{चयन} (cayan) — a selection; a choice (formal)}

\label{lesson:HI-C550-cayan}



\begin{quote}
Before the new one: say the Hindi for a ban, a prohibition, then the Hindi for favourite.
\end{quote}

\subsection*{You'll want to know: \hi{चयन}}
\textbf{\hi{चयन}} — \emph{cayan} — \textquotedblleft{}a selection; a choice (formal)\textquotedblright{}.

\begin{quote}
\hi{टीम} \hi{के} \hi{लिए} \hi{ग्यारह} \hi{खिलाड़ियों} \hi{का} \hi{चयन} \hi{हुआ।}
\end{quote}

\begin{quote}
\emph{Ṭīm ke liye gyārah khilāṛiyoṁ kā cayan huā.}
\end{quote}

\begin{quote}
Eleven players were picked for the team.
\end{quote}

Masculine.

A formal word.

\begin{cousinweb}
\textbf{\hi{चयन}} (\emph{cayan}) is a tatsama, taken whole from Sanskrit \emph{cayana}, \textquotedblleft{}piling up, collecting\textquotedblright{} (Monier-Williams). Hindi uses it for picking out, a selection.
\end{cousinweb}

\begin{grammarlens}[title={What you've built}]
One more word for saying what happens if.
\end{grammarlens}

\subsection*{Your turn}
\begin{itemize}
\raggedright
  \item \emph{Say it:} \emph{cayan}
  \item \emph{Say it:} \emph{cayan}, once more
  \item \emph{Say it:} \emph{Ṭīm ke liye gyārah khilāṛiyoṁ kā cayan huā.} — Eleven players were picked for the team.
  \item \emph{Recall:} say the Hindi for a ban, a prohibition, then the Hindi for favourite, then say \emph{cayan} again
\end{itemize}

\begin{tcolorbox}[breakable,colback=teal!4,colframe=teal!35!black,title={Before you move on}]
What does \hi{चयन} mean? (\textquotedblleft{}A selection; a choice (formal)\textquotedblright{}.) Say it once more.
\end{tcolorbox}
